% ID: SEL-D0023-OA-E02
% ESTADO: validada. S024, cotejada con el examen.
\textbf{a)} Aplicamos producto y cadena a $f$. Para $g$,
escribimos $g(x)=\ln(3x)e^{-2x}$ y derivamos para $x>0$.
\[
\begin{aligned}
f'(x)&=2x(3x^3+5x)^3+3(x^2-1)(3x^3+5x)^2(9x^2+5),\\
g'(x)&=e^{-2x}\bigl(1/x-2\ln(3x)\bigr).
\end{aligned}
\]
\textbf{b)} Por el cociente, $h'(x)=-9/(2x+1)^2$.
En $x=1$, $h(1)=3$ y $h'(1)=-1$.
La tangente es $y-3=-(x-1)$, es decir, $\boxed{y=4-x}$.

\textbf{c)} Al escribir $h(x)=3/2+9/[2(2x+1)]$ se ve que
$h$ tiende a $3/2$ cuando $x\to\pm\infty$, y que diverge con signos
opuestos a ambos lados de $-1/2$. Las asíntotas son
$\boxed{x=-1/2}$ y $\boxed{y=3/2}$.
